% GENERATED FILE. Edit canonical lessons, then run npm run generate:books.
% canonical-source-hash: fnv1a64:4e0d3f016c0f3a7e
% canonical-lessons: GE-C14-haben, GE-C14-hab-stamm, GE-C14-du-hast, GE-C14-er-hat, GE-C14-b-verlust, GE-C14-einen, GE-C14-habere, GE-C14-practice

\chapter{To Have}
\label{ch:haben}
\hlchaptermodality{18}

\begin{chapteropening}
\textbf{By the end of this chapter:} \emph{I can use every present form of German's verb for “to have”, building the four the ending machine already gives me and knowing the two it gets wrong.}

haben is the weak ending machine on the stem hab-, and the chapter is organised around where that machine fails: it predicts four of the six forms correctly, so those get one lesson that makes the reader build them rather than read them, and du hast and er hat get a lesson each because nothing generates them. The b-drop, the einen that Ich habe einen Bruder needs, and the habēre false cognate follow; the grid is printed only in the practice lesson.
\end{chapteropening}

\section[haben]{haben — \textquotedblleft{}to have\textquotedblright{}}

\label{lesson:GE-C14-haben}



\begin{quote}
German's workhorse verb, and the twin of a word you have said every day of your life.
\end{quote}

\subsection*{You'll want to know: haben}
\begin{quote}
\textbf{haben} — \textquotedblleft{}to have.\textquotedblright{}
\end{quote}

Same word as English \emph{have}, worn differently. Chop the \emph{-en} the way you did with \emph{machen} and the stem is \emph{hab-}, which is where every form in this chapter starts.

Two syllables, and the first one carries the weight: \emph{HAH-ben}.

\begin{sounds}
\begin{itemize}
\raggedright
  \item \texttt{long-a} — the \emph{a} of \emph{haben} is long and open, the vowel of English \emph{father}, not the clipped \emph{a} of \emph{hat}.
  \item \texttt{d-t-final} — that long \emph{a} does \textbf{not} survive into every form. Two of them coming up are short and clipped, and both end in a crisp, fully released \emph{t}.
\end{itemize}
\end{sounds}

\subsection*{Your turn}
Say these aloud:

\begin{itemize}
\raggedright
  \item \textquotedblleft{}haben\textquotedblright{} — long \emph{a}, soft second syllable
  \item the stem on its own — \textquotedblleft{}hab-\textquotedblright{}
  \item the twins — \textquotedblleft{}haben, have\textquotedblright{}
\end{itemize}

\emph{Twice through:} \textquotedblleft{}haben, haben — to have.\textquotedblright{}

\begin{tcolorbox}[breakable,colback=teal!4,colframe=teal!35!black,title={Before you move on}]
Say \textquotedblleft{}to have.\textquotedblright{} (\textbf{Haben}.) What is its stem? (\emph{\textbf{Hab-}}.) Is the \emph{a} long or short? (\textbf{Long} — as in \emph{father}.) Which English verb is it? (\emph{\textbf{Have}}.) Next: the stem through the ending machine you already own.
\end{tcolorbox}
\section[ich habe, wir haben, ihr habt]{hab- through the machine}

\label{lesson:GE-C14-hab-stamm}



\begin{quote}
You own an ending machine: a stem, and \emph{-e}, \emph{-st}, \emph{-t}, \emph{-en}. Run \emph{hab-} through it before reading on, and keep what you get.
\end{quote}

\begin{grammarlens}[title={Grammar Lens: four forms you already knew}]
\begin{quote}
On \emph{hab-}, the machine is \textbf{right four times out of six}.
\end{quote}

\emph{Ich} takes \emph{-e}, so \emph{ich} \textbf{habe}. \emph{Wir} takes \emph{-en}, so \emph{wir} \textbf{haben} — which is the infinitive over again, the way the \emph{wir} form goes throughout German. \emph{Ihr} takes \emph{-t}, so \emph{ihr} \textbf{habt}. And \emph{sie} takes \emph{-en} like \emph{wir}, so \emph{sie} \textbf{haben}.

Nothing there is new. Those four are the rule from your first verbs applied to one more stem — so do not read them off the page. Take each ending, put it on \emph{hab-}, and hear that you already knew the answer.

The chapter is about the other two.
\end{grammarlens}

\subsection*{Your turn}
Build each one before you say it — take the ending first, then put it on \emph{hab-}.

Say these aloud:

\begin{itemize}
\raggedright
  \item the ending \emph{ich} takes, then the form — \textquotedblleft{}-e … ich habe\textquotedblright{}
  \item the ending \emph{wir} takes, then the form — \textquotedblleft{}-en … wir haben\textquotedblright{}
  \item the ending \emph{ihr} takes, then the form — \textquotedblleft{}-t … ihr habt\textquotedblright{}
  \item \emph{sie} borrows \emph{wir}'s ending — \textquotedblleft{}sie haben\textquotedblright{}
  \item now prove the machine is not about this verb — run \emph{wohn-} through it: \textquotedblleft{}ich wohne, wir wohnen, ihr wohnt\textquotedblright{}
\end{itemize}

\emph{Twice through:} \textquotedblleft{}ich habe, wir haben, ihr habt, sie haben.\textquotedblright{}

\begin{tcolorbox}[breakable,colback=teal!4,colframe=teal!35!black,title={Before you move on}]
Build \textquotedblleft{}we have\textquotedblright{} without being given it. (\textbf{Wir haben}.) Build \textquotedblleft{}you all have.\textquotedblright{} (\textbf{Ihr habt}.) Which other form does \emph{wir haben} match? (\textbf{The infinitive}.) Did you have to learn any of those four, or work them out? (\textbf{Work them out} — the endings were already yours.) How many of the six does the rule get right? (\textbf{Four}.) Next: the first one it gets wrong.
\end{tcolorbox}
\section[du hast]{du hast — \textquotedblleft{}you have\textquotedblright{}}

\label{lesson:GE-C14-du-hast}



\begin{quote}
The machine says \emph{du} plus \emph{-st} on the stem \emph{hab-}. Say what that would give, and then hear what German actually says.
\end{quote}

\subsection*{You'll want to know: du hast}
\begin{quote}
\textbf{du hast} — \textquotedblleft{}you have,\textquotedblright{} to one person you are on familiar terms with.
\end{quote}

The rule predicts \emph{hab-st}. German says \textbf{hast}. The \emph{b} is simply gone, and the long \emph{a} goes with it: the vowel is short and clipped now, and the \emph{-st} lands crisply on the end.

This is not a spelling convenience. It is one of exactly two places in the whole verb where the machine fails, and it is worth meeting on its own so that it is a form you know rather than a form you derive wrongly.

\begin{sounds}
\begin{itemize}
\raggedright
  \item \texttt{vowel-a} — short and clipped, the \emph{a} of English \emph{hat}, not the long \emph{a} of \emph{haben}. The form is one syllable: \emph{hast}, rhyming with English \emph{mast}.
  \item \texttt{d-t-final} — the \emph{-st} is fully released; German does not soften it.
\end{itemize}
\end{sounds}

\subsection*{Your turn}
Say these aloud:

\begin{itemize}
\raggedright
  \item \textquotedblleft{}du hast\textquotedblright{}
  \item the long against the short — \textquotedblleft{}ich habe, du hast\textquotedblright{}
  \item what the rule wanted, then what German says — \textquotedblleft{}hab-st … hast\textquotedblright{}
\end{itemize}

\emph{Twice through:} \textquotedblleft{}du hast, du hast.\textquotedblright{}

\begin{tcolorbox}[breakable,colback=teal!4,colframe=teal!35!black,title={Before you move on}]
Say \textquotedblleft{}you have,\textquotedblright{} informally. (\textbf{Du hast}.) Which letter went missing? (\textbf{The b}.) What happened to the vowel? (\textbf{It shortened}.) What would the rule have predicted? (\emph{\textbf{Hab-st}}.) Next: the second place the machine fails, and it fails the same way.
\end{tcolorbox}
\section[er hat]{er hat — \textquotedblleft{}he has\textquotedblright{}}

\label{lesson:GE-C14-er-hat}



\begin{quote}
One more break, and then the verb behaves for good. The machine says \emph{-t} on \emph{hab-}.
\end{quote}

\subsection*{You'll want to know: er hat}
\begin{quote}
\textbf{er hat} — \textquotedblleft{}he has.\textquotedblright{} Also \emph{sie hat}, \textquotedblleft{}she has,\textquotedblright{} and \emph{es hat}, \textquotedblleft{}it has.\textquotedblright{}
\end{quote}

The rule predicts \emph{hab-t}. German says \textbf{hat}: the same \emph{b} gone, the same vowel shortened. One syllable, rhyming with English \emph{hut} rather than \emph{hat} — the German \emph{a} sits further back than the English one.

One form covers \emph{er}, \emph{sie} and \emph{es}, exactly as the third-person \emph{-t} always does. So the two breaks in this verb are neighbours: \emph{du} \textbf{hast}, \emph{er} \textbf{hat}, and every other form is the machine.

\subsection*{Your turn}
Say these aloud:

\begin{itemize}
\raggedright
  \item \textquotedblleft{}er hat, sie hat, es hat\textquotedblright{}
  \item the two breaks together — \textquotedblleft{}du hast, er hat\textquotedblright{}
  \item the whole shape — \textquotedblleft{}habe … hast … hat … haben\textquotedblright{}
\end{itemize}

\emph{Twice through:} \textquotedblleft{}du hast, er hat.\textquotedblright{}

\begin{tcolorbox}[breakable,colback=teal!4,colframe=teal!35!black,title={Before you move on}]
Say \textquotedblleft{}he has.\textquotedblright{} (\textbf{Er hat}.) Say \textquotedblleft{}she has.\textquotedblright{} (\textbf{Sie hat}.) Which two forms of this verb break the rule? (\emph{\textbf{Hast}} \textbf{and} \emph{\textbf{hat}}.) What do both of them drop? (\textbf{The b}.) Next: why they drop it, and the English verb that does the identical thing.
\end{tcolorbox}
\section[der verlorene b]{The b that fell out}

\label{lesson:GE-C14-b-verlust}



\begin{quote}
\emph{Hast} and \emph{hat} both lost a \emph{b}. That is not two accidents.
\end{quote}

\begin{grammarlens}[title={Grammar Lens: one shortcut, in exactly two cells}]
\begin{quote}
\emph{Haben} keeps its \textbf{b} everywhere except before a \emph{-st} or a \emph{-t} ending, where the \emph{b} drops and the vowel shortens with it.
\end{quote}

Say \emph{hab-st} and \emph{hab-t} at speed and you can feel why: a \emph{b} squeezed between a vowel and a hard \emph{t} is work, and centuries of speakers declined to do it. The consonant wore away and the two shortest forms are what is left.

Now listen to English, which is the same verb:

\noindent
\begin{tabularx}{\linewidth}{@{}>{\raggedright\arraybackslash}X>{\raggedright\arraybackslash}X>{\raggedright\arraybackslash}X@{}}
\toprule
\textbf{person} & \textbf{English} & \textbf{German} \\
\midrule
infinitive & \emph{have} & \emph{haben} \\
second person, older & \emph{thou ha\textbf{st}} & \emph{du ha\textbf{st}} \\
third person & \emph{he ha\textbf{s}} & \emph{er ha\textbf{t}} \\
\bottomrule
\end{tabularx}

English lost the same consonant, in the same two cells, of the same verb — and kept the archaic \emph{hast} long enough for us to see it. Two languages, one shortcut, inherited together rather than invented twice.
\end{grammarlens}

\subsection*{Your turn}
Say these aloud:

\begin{itemize}
\raggedright
  \item where the b survives — \textquotedblleft{}habe, haben, habt\textquotedblright{}
  \item where it does not — \textquotedblleft{}hast, hat\textquotedblright{}
  \item the English echo — \textquotedblleft{}thou hast, du hast; he has, er hat\textquotedblright{}
\end{itemize}

\emph{Twice through:} \textquotedblleft{}hast, hat — no b.\textquotedblright{}

\begin{tcolorbox}[breakable,colback=teal!4,colframe=teal!35!black,title={Before you move on}]
In which two forms does the \emph{b} drop? (\emph{\textbf{Hast}} \textbf{and} \emph{\textbf{hat}}.) What else changes with it? (\textbf{The vowel shortens}.) Which English form still shows the same loss? (\emph{\textbf{He has}} — or archaic \emph{\textbf{thou hast}}.) Next: something to actually have.
\end{tcolorbox}
\section[einen]{einen — \textquotedblleft{}a,\textquotedblright{} when it is the thing you have}

\label{lesson:GE-C14-einen}



\begin{quote}
\emph{Ich habe …} wants something after it. The commonest thing to put there is \textquotedblleft{}a something,\textquotedblright{} and German asks one small question first.
\end{quote}

\subsection*{You'll want to know: einen}
\begin{quote}
\textbf{einen} — \textquotedblleft{}a,\textquotedblright{} for a \textbf{der}-word that is the thing you have.
\end{quote}

\emph{Ein} is the number \emph{one} doing a second job, which is exactly what English \emph{a} is — a worn-down \emph{one}. German keeps the resemblance visible: \emph{ein}, \emph{eins}.

The \emph{-en} on the end is the part English threw away. German marks a \emph{der}-word differently when it is the thing being had rather than the thing doing the having, and \emph{einen} is that marked form:

\begin{quote}
\textbf{Ich habe einen Bruder.} — \textquotedblleft{}I have a brother.\textquotedblright{}
\end{quote}

\emph{Der Bruder} is a \emph{der}-word, and he is what you have, so \emph{ein} becomes \emph{einen}. Nothing else in the sentence moves.

Take this one as a word for now, not as a system. The system behind it — German's cases — gets a chapter of its own later; this is the first place the book cannot avoid it, because a sentence with \emph{haben} in it needs something to have.

\subsection*{Your turn}
Say these aloud:

\begin{itemize}
\raggedright
  \item \textquotedblleft{}einen\textquotedblright{}
  \item \textquotedblleft{}Ich habe einen Bruder.\textquotedblright{}
  \item it with the two broken forms — \textquotedblleft{}du hast einen Bruder; er hat einen Bruder\textquotedblright{}
\end{itemize}

\emph{Twice through:} \textquotedblleft{}Ich habe einen Bruder.\textquotedblright{}

\begin{tcolorbox}[breakable,colback=teal!4,colframe=teal!35!black,title={Before you move on}]
Say \textquotedblleft{}I have a brother.\textquotedblright{} (\textbf{Ich habe einen Bruder}.) Which word did \emph{ein} become? (\emph{\textbf{Einen}}.) Why? (\textbf{Because \emph{der Bruder} is the thing being had}.) Which English word is \emph{ein} underneath? (\emph{\textbf{One}}.) Next: the trap this verb is famous for.
\end{tcolorbox}
\section[haben ist nicht habēre]{haben is not habēre}

\label{lesson:GE-C14-habere}



\begin{quote}
German \emph{haben}, Latin \emph{habēre}. Same consonants, same meaning, and the best false cognate in the business.
\end{quote}

\begin{cousinweb}
Two verbs that look like one word:

\noindent
\begin{tabularx}{\linewidth}{@{}>{\raggedright\arraybackslash}X>{\raggedright\arraybackslash}X>{\raggedright\arraybackslash}X@{}}
\toprule
\textbf{language} & \textbf{verb} & \textbf{meaning} \\
\midrule
German & \textbf{haben} & to have \\
Latin & \textbf{habēre} & to have \\
\bottomrule
\end{tabularx}

\textbf{They are not related.} It is a coincidence, and one of the most famous in linguistics. Follow each one back and they go in opposite directions.

\emph{Haben} comes from Germanic *\emph{habjaną}, from Proto-Indo-European *\emph{kap-}, \textquotedblleft{}\textbf{to seize}.\textquotedblright{} That root's Latin child is \emph{capere} — which is where English gets \textbf{capture}, \textbf{captive}, \textbf{capable} and \textbf{accept}.

\emph{Habēre} comes from a different root entirely, *\emph{g\textsuperscript{h}ab\textsuperscript{h}-}, and its English descendant is … \textbf{give}.

So German \emph{haben} is a cousin of \emph{capture}, and Latin \emph{habēre} — the source of French \emph{avoir} and Italian \emph{avere} — is a cousin of \emph{give}. Seizing and giving, which is about as far apart as two verbs get.
\end{cousinweb}

\begin{culture}
This is worth carrying beyond one verb. Resemblance is not evidence. Two words can match in sound and sense and still have nothing to do with each other, and the only way to tell is to follow both back rather than to compare them where they stand.

Most of this book's cousin pairs are real. This one is the control that proves the others were checked.
\end{culture}

\subsection*{Your turn}
Say these aloud:

\begin{itemize}
\raggedright
  \item the pair that is not a pair — \textquotedblleft{}haben, habēre\textquotedblright{}
  \item where German's goes — \textquotedblleft{}haben, kap-, capture\textquotedblright{}
  \item where Latin's goes — \textquotedblleft{}habēre, g\textsuperscript{h}ab\textsuperscript{h}-, give\textquotedblright{}
\end{itemize}

\emph{Twice through:} \textquotedblleft{}haben is not habēre.\textquotedblright{}

\begin{tcolorbox}[breakable,colback=teal!4,colframe=teal!35!black,title={Before you move on}]
Are \emph{haben} and \emph{habēre} related? (\textbf{No} — a coincidence.) Which English word is \emph{haben}'s cousin? (\emph{\textbf{Capture}}, from \emph{capere}.) And \emph{habēre}'s? (\emph{\textbf{Give}}.) What did \emph{haben}'s root originally mean? (\textquotedblleft{}\textbf{To seize}.\textquotedblright{}) Next: the whole verb, said through.
\end{tcolorbox}
\section[Practice]{Practice — the whole of haben}

\label{lesson:GE-C14-practice}



\begin{quote}
No new words. Four of the six forms below cost you nothing, and the other two have each had a lesson, so this is the first place it is safe to set them out together.
\end{quote}

\begin{grammarlens}[title={Grammar Lens: the six, assembled}]
\noindent
\begin{tabularx}{\linewidth}{@{}>{\raggedright\arraybackslash}X>{\raggedright\arraybackslash}X@{}}
\toprule
\textbf{singular} & \textbf{plural} \\
\midrule
ich \textbf{habe} & wir \textbf{haben} \\
du \textbf{hast} & ihr \textbf{habt} \\
er/sie/es \textbf{hat} & sie/Sie \textbf{haben} \\
\bottomrule
\end{tabularx}

Run it aloud and hear where the \emph{b} goes missing: \emph{habe} … \textbf{hast} … \textbf{hat} … \emph{haben}. Four forms carry it, two do not, and the two that do not are the two shortest.
\end{grammarlens}

\subsection*{The exchange}
\noindent
\begin{tabularx}{\linewidth}{@{}>{\raggedright\arraybackslash}X>{\raggedright\arraybackslash}X@{}}
\toprule
\textbf{German} & \textbf{English} \\
\midrule
\emph{Hast du einen Bruder?} & Do you have a brother? \\
\emph{Und du?} & And you? \\
\emph{Ich habe eine Schwester.} & I have a sister. \\
\emph{Wir haben Brot und Wein.} & We have bread and wine. \\
\emph{Gut! Dann haben wir alles.} & Good! Then we have everything. \\
\bottomrule
\end{tabularx}

Two things slip in without lessons of their own, because the translation carries them: \textbf{eine}, which is what \emph{ein} becomes for a \emph{die}-word, and \textbf{alles} (\textquotedblleft{}everything\textquotedblright{}). And notice \emph{Hast du …?} — a question puts the verb first, which is the one move German makes that English needs \emph{do} for.

\subsection*{Your turn}
Say these aloud:

\begin{itemize}
\raggedright
  \item the six in order — \textquotedblleft{}habe, hast, hat, haben, habt, haben\textquotedblright{}
  \item the two breaks alone — \textquotedblleft{}hast, hat\textquotedblright{}
  \item \textquotedblleft{}Ich habe einen Bruder.\textquotedblright{} \textquotedblleft{}Wir haben Brot und Wein.\textquotedblright{}
  \item the whole exchange, both voices
\end{itemize}

\emph{Twice through:} \textquotedblleft{}habe, hast, hat, haben, habt, haben.\textquotedblright{}

\begin{tcolorbox}[breakable,colback=teal!4,colframe=teal!35!black,title={Before you move on}]
Run the six without looking. Which two drop the \emph{b}? (\emph{\textbf{Hast}} \textbf{and} \emph{\textbf{hat}}.) Say \textquotedblleft{}I have a brother.\textquotedblright{} (\textbf{Ich habe einen Bruder}.) Is \emph{haben} related to Latin \emph{habēre}? (\textbf{No}.) Which English word is it related to? (\emph{\textbf{Capture}}.) Next: the one everyday sentence where German refuses this verb altogether.
\end{tcolorbox}
